\documentclass[a4paper,12pt]{article}
\usepackage[T2A]{fontenc}
\usepackage[utf8]{inputenc}
\usepackage[russian]{babel}
\usepackage{amsmath,amssymb,mathtools}
\usepackage{helvet}
\renewcommand{\familydefault}{\sfdefault}
\usepackage{icomma}
\usepackage[a4paper,left=19mm,right=19mm,top=18mm,bottom=20mm]{geometry}
\usepackage{microtype}
\usepackage{tikz}
\usetikzlibrary{calc,arrows.meta,positioning,shapes.geometric}
\usepackage{tkz-euclide}
\usepackage{pgfplots}
\pgfplotsset{compat=1.18}
\usepackage{tabularx,booktabs}
\usepackage{enumitem}
\usepackage{fancyhdr}
\usepackage{setspace}
\usepackage{xcolor}
\usepackage{hyperref}
\usepackage{parskip}
\usepackage{needspace}
\definecolor{accent}{HTML}{176B73}
\definecolor{darkaccent}{HTML}{124A50}
\definecolor{softgray}{HTML}{66767B}
\definecolor{textgray}{HTML}{2B3438}
\definecolor{warm}{HTML}{B66A3C}
\hypersetup{colorlinks=true,linkcolor=darkaccent,urlcolor=darkaccent}
\setlength{\parindent}{0pt}
\setlength{\parskip}{4.5pt}
\setstretch{1.22}
\setlist[itemize]{leftmargin=7mm,itemsep=2pt,topsep=3pt}
\setlist[enumerate]{leftmargin=8mm,itemsep=5pt,topsep=4pt}
\pagestyle{fancy}
\fancyhf{}
\fancyfoot[C]{\small\color{softgray}\thepage}
\renewcommand{\headrulewidth}{0pt}
\newcommand{\sect}[1]{\vspace{2mm}{\Large\bfseries\color{darkaccent}#1}\par\vspace{1mm}{\color{accent!50}\rule{\linewidth}{.55pt}}\par\vspace{1mm}}
\newcommand{\idea}[1]{\textbf{\color{darkaccent}#1}}
\newcommand{\checkline}[1]{\noindent$\square$\enspace #1\par}
\begin{document}
\thispagestyle{empty}
\begin{center}
\vspace*{33mm}
{\small\color{softgray}МАТЕМАТИКА \textbullet{} ОГЭ \textbullet{} ПОВТОРЕНИЕ И УГЛУБЛЕНИЕ}\par
\vspace{11mm}
{\fontsize{29}{32}\selectfont\bfseries\color{darkaccent}Чек-лист}\par
\vspace{6mm}
{\LARGE\bfseries Графики, неравенства и геометрические доказательства}\par
\vspace{8mm}
{\large Модуль и выколотые точки \textbullet{} Смеси\\Биссектрисы трапеции \textbullet{} Высоты и окружность}\par
\vspace{23mm}
{\large София Кальней}\par
\vspace{4mm}
{\normalsize 11 октября 2026 г.}\par
\vfill
{\small\color{textgray}Лёвин Артём Александрович эксклюзивно для Софии Кальней}\par
\vspace{20mm}
\end{center}
\begin{tikzpicture}[remember picture,overlay]
\draw[accent!50,line width=1.15pt]
 ([xshift=8mm,yshift=7mm]current page.south west)
 .. controls +(47mm,21mm) and +(-56mm,-5mm) ..
 ([xshift=-8mm,yshift=14mm]current page.south east);
\draw[warm!40,line width=.65pt]
 ([xshift=16mm,yshift=5mm]current page.south west)
 .. controls +(55mm,17mm) and +(-66mm,-8mm) ..
 ([xshift=-16mm,yshift=11mm]current page.south east);
\end{tikzpicture}
\clearpage

\sect{1. Модуль и преобразования графиков}
\checkline{Различаю действия $|f(x)|$, $-|f(x)|$ и $f(|x|)$.}
\[
\begin{aligned}
 y=|f(x)|&:\quad\text{части ниже оси }Ox\text{ отражаются вверх};\\
 y=-|f(x)|&:\quad\text{весь график }|f(x)|\text{ отражается вниз};\\
 y=f(|x|)&:\quad\text{правую часть графика }f(x)\text{ отражаем влево}.
\end{aligned}
\]
\idea{Пример.} Для $y=-|x^2-x-2|$ сначала строим $y=x^2-x-2$, затем $y=|x^2-x-2|$, и только после этого отражаем результат относительно $Ox$. Нули: $x=-1$ и $x=2$; весь итоговый график расположен в полуплоскости $y\le0$.

\idea{Самопроверка.} Знак минус \emph{перед всем модулем} меняет знак всех значений функции. Модуль в аргументе $f(|x|)$ действует иначе.

\sect{2. Дробь с модулем: ОДЗ, сокращение и график}
\checkline{Определяю ОДЗ \emph{до} сокращения дроби.}
\checkline{Отличаю выколотую точку от асимптоты.}

\idea{Пример с занятия.} Построим график функции
\[
g(x)=\frac{\frac32|x|-1}{|x|-\frac32x^2}.
\]
Так как $x^2=|x|^2$, то
\[
|x|-\frac32x^2=|x|\left(1-\frac32|x|\right).
\]
\textbf{ОДЗ:} $x\ne0$, $x\ne-\frac23$, $x\ne\frac23$. При этих ограничениях
\[
g(x)=\frac{-\left(1-\frac32|x|\right)}{|x|\left(1-\frac32|x|\right)}
=\boxed{-\frac1{|x|}}.
\]
Сокращение не возвращает исключённые значения. При $x=\pm\frac23$ на графике имеются \textbf{выколотые точки} с ординатой $-\frac32$; при $x=0$ --- вертикальная асимптота.

\begin{center}
\begin{tikzpicture}
\begin{axis}[
 width=11.8cm,height=7cm,
 axis lines=middle,axis line style={->,textgray},
 xmin=-3.2,xmax=3.2,ymin=-4.4,ymax=.45,
 xtick={-3,-2,-1,1,2,3},ytick={-4,-3,-2,-1},
 xlabel={$x$},ylabel={$y$},
 x label style={at={(axis description cs:1.02,.91)}},
 y label style={at={(axis description cs:.52,1.02)}},
 tick label style={font=\scriptsize},
 grid=major,grid style={gray!15},
 clip=false]
\addplot[accent,thick,domain=-3.1:-.22,samples=170,restrict y to domain=-4.4:.45] {-1/abs(x)};
\addplot[accent,thick,domain=.22:3.1,samples=170,restrict y to domain=-4.4:.45] {-1/abs(x)};
\addplot[only marks,mark=o,mark size=3pt,mark options={draw=warm,line width=1.2pt,fill=white}] coordinates {(-.666667,-1.5) (.666667,-1.5)};
\node[font=\scriptsize,anchor=north east,fill=white,inner sep=1pt] at (axis cs:-.67,-1.5) {$(-\frac23;-\frac32)$};
\node[font=\scriptsize,anchor=north west,fill=white,inner sep=1pt] at (axis cs:.67,-1.5) {$(\frac23;-\frac32)$};
\end{axis}
\end{tikzpicture}
\end{center}
\textbf{Последовательность построения:} $y=1/x\ \to\ y=1/|x|\ \to\ y=-1/|x|$; затем исключаем точки, запрещённые первоначальным знаменателем.

\sect{3. Прямая $y=kx$ и число общих точек}
\checkline{Помню, что прямая $y=kx$ проходит через начало координат.}
\checkline{Учитываю исключённые точки исходной функции.}

Для графика $g(x)$ из предыдущего раздела ищем значения $k$, при которых прямая $y=kx$ не имеет с ним общих точек.

Если $x>0$, то $-1/x=kx$, откуда $k=-1/x^2<0$. Единственный \emph{пропавший} случай соответствует $x=2/3$:
\[
k=-\frac1{(2/3)^2}=-\frac94.
\]
Если $x<0$, то $1/x=kx$, откуда $k=1/x^2>0$. Исключённое значение $x=-2/3$ даёт $k=9/4$. При $k=0$ прямая является осью $Ox$ и также не пересекает график, лежащий ниже оси.
\[
\boxed{k\in\left\{-\frac94,\ 0,\ \frac94\right\}.}
\]
\idea{Важное различие.} Выколотая точка может убрать единственное пересечение с прямой. Поэтому исследование упрощённой формулы без ОДЗ даёт неполный ответ.

\sect{4. Дробно-рациональные неравенства}
\checkline{Нахожу нули числителя и запрещённые нули знаменателя.}
\checkline{Проверяю знаки на промежутках, затем выбираю подходящие.}

\idea{Пример с занятия.} Решим
\[
\frac{-10}{(x-3)^2-5}\ge0.
\]
Числитель отрицателен и никогда не обращается в нуль. Чтобы дробь была неотрицательной, знаменатель обязан быть \emph{строго} отрицательным:
\[
(x-3)^2-5<0\quad\Longleftrightarrow\quad
-\sqrt5<x-3<\sqrt5.
\]
\[
\boxed{x\in(3-\sqrt5;\ 3+\sqrt5).}
\]
Обе границы исключены: там знаменатель равен нулю, даже несмотря на знак $\ge$ в условии.

\idea{Универсальный метод интервалов.} Представьте выражение одной дробью, определите ОДЗ, отметьте критические точки, проверьте знак дроби на каждом промежутке. Корни числителя включаются только при нестрогом неравенстве и допустимом знаменателе.

\sect{5. Текстовая задача на смеси}
\checkline{Различаю массу раствора, концентрацию и массу вещества.}
\checkline{Для каждого способа смешивания записываю своё уравнение.}

Пусть $x$ и $y$ --- доли кислоты в первом и втором растворах. В первом сосуде $40$ кг, во втором $20$ кг. Общая концентрация при полном смешивании равна $33\%$, а при смешивании равных масс --- $47\%$.
\[
\underbrace{40x+20y}_{\text{масса кислоты}}=60\cdot0{,}33,
\qquad \frac{x+y}{2}=0{,}47.
\]
\[
\begin{cases}2x+y=0{,}99,\\x+y=0{,}94\end{cases}
\quad\Longrightarrow\quad x=0{,}05,\quad y=0{,}89.
\]
\textbf{Ответ:} в первом растворе $5\%$ кислоты. Проверка: $40\cdot0{,}05+20\cdot0{,}89=19{,}8$ кг, что составляет $33\%$ от $60$ кг.

\idea{Частая ошибка.} Среднее арифметическое концентраций допустимо при смешивании \emph{равных масс}. При разных массах нужна взвешенная сумма.

\sect{6. Биссектрисы углов при боковой стороне трапеции}
\checkline{Использую сумму односторонних углов при параллельных основаниях.}
\checkline{После доказательства прямого угла применяю Пифагора.}

Пусть в трапеции $ABCD$ основания $AD\parallel BC$, а $AB$ --- боковая сторона. Биссектрисы углов $DAB$ и $ABC$ пересекаются в точке $K$.
\[
\angle DAB+\angle ABC=180^\circ.
\]
Поскольку $AK$ и $BK$ делят эти углы пополам,
\[
\angle KAB+\angle ABK=90^\circ
\quad\Longrightarrow\quad
\boxed{\angle AKB=90^\circ}.
\]
\textbf{Пример с занятия:} если $AK=18$, $BK=24$, то
\[
AB=\sqrt{18^2+24^2}=\boxed{30}.
\]
\idea{Другой приём с биссектрисами.} Если биссектрисы углов $A$ и $D$ пересекаются в точке $M$ на стороне $BC$, то расстояния от $M$ до прямых $AB$, $AD$, $CD$ равны. Каждая биссектриса задаёт равенство расстояний до \emph{сторон своего угла}.

\Needspace{150mm}
\sect{7. Две высоты и окружность с диаметром $BC$}
\checkline{В тупоугольном треугольнике продлеваю стороны при необходимости.}
\checkline{Распознаю четыре точки на одной окружности.}

В $\triangle ABC$ угол $BAC$ тупой. Проведены высоты $BB_1\perp AC$ и $CC_1\perp AB$ (основания лежат на продолжениях сторон). Надо доказать:
\[
\boxed{\triangle AB_1C_1\sim\triangle ABC}.
\]
\begin{center}
\begin{tikzpicture}[scale=.82]
\tkzDefPoint(0,0){A}
\tkzDefPoint(5,0){B}
\tkzDefPoint(-2,3){C}
\tkzDefPoint(1.53846,-2.30769){B1}
\tkzDefPoint(-2,0){C1}
\tkzDrawSegments[color=softgray,dashed](C,B1 B,C1)
\tkzDrawPolygon[color=textgray,line width=1pt](A,B,C)
\tkzDrawSegments[color=accent,line width=1pt](B,B1 C,C1 B1,C1)
\tkzDrawSegments[color=softgray,dashed](A,B1 A,C1)
\tkzMarkRightAngle[size=.23](A,B1,B)
\tkzMarkRightAngle[size=.23](A,C1,C)
\tkzDrawPoints[size=2.2,fill=darkaccent](A,B,C,B1,C1)
\tkzLabelPoint[above right](A){$A$}
\tkzLabelPoint[right](B){$B$}
\tkzLabelPoint[above](C){$C$}
\tkzLabelPoint[below right](B1){$B_1$}
\tkzLabelPoint[left](C1){$C_1$}
\end{tikzpicture}
\end{center}
\textbf{Доказательство.} Углы $BB_1C$ и $BC_1C$ прямые, следовательно, точки $B,B_1,C,C_1$ лежат на одной окружности с диаметром $BC$. Вписанные углы $CB_1C_1$ и $CBC_1$ опираются на хорду $CC_1$ и равны. Поэтому
\[
\angle AB_1C_1=\angle ABC.
\]
Кроме того, $AB_1$ и $AC_1$ являются продолжениями лучей $AC$ и $AB$ в противоположные стороны, поэтому $\angle B_1AC_1=\angle BAC$. Подобие следует по двум углам.

\idea{Запомните.} Две высоты одного треугольника часто указывают на окружность, построенную по общей гипотенузе прямоугольных треугольников.

\clearpage
\begin{center}
{\Large\bfseries\color{darkaccent}Алгоритм: дробная функция с модулем}\par\vspace{6mm}
{\small\color{softgray}Как сохранить все ограничения после преобразований}
\end{center}
\begin{center}
\begin{tikzpicture}[
 >=Stealth, node distance=10mm,
 start/.style={ellipse,draw=darkaccent,line width=.9pt,align=center,text width=8.0cm,minimum height=10mm,fill=accent!5},
 flowstep/.style={rounded corners=2mm,rectangle,draw=accent!75,line width=.8pt,align=center,text width=10.7cm,minimum height=15mm,inner sep=4mm},
 arrow/.style={->,line width=1.1pt,draw=darkaccent}
]
\node[start] (a) {Исходная дробная функция $y=\frac{P(x)}{Q(x)}$};
\node[flowstep,below=of a] (b) {\textbf{1. ОДЗ}: решить $Q(x)\ne0$\\и записать все исключённые значения};
\node[flowstep,below=of b] (c) {\textbf{2. Преобразовать}: использовать $x^2=|x|^2$,\\разложить числитель и знаменатель на множители};
\node[flowstep,below=of c] (d) {\textbf{3. Сократить общие множители},\\сохраняя первоначальную ОДЗ};
\node[flowstep,below=of d] (e) {\textbf{4. Построить график} упрощённой функции,\\учитывая модуль и отражения};
\node[flowstep,below=of e] (f) {\textbf{5. Отметить выколотые точки}\\и асимптоты, если они возникают};
\node[flowstep,below=of f] (g) {\textbf{6. Для пересечений} составить уравнение\\с учётом ОДЗ и проверить найденные корни};
\node[start,below=of g] (z) {Готовый график и корректный ответ};
\draw[arrow] (a)--(b);
\draw[arrow] (b)--(c);
\draw[arrow] (c)--(d);
\draw[arrow] (d)--(e);
\draw[arrow] (e)--(f);
\draw[arrow] (f)--(g);
\draw[arrow] (g)--(z);
\end{tikzpicture}
\end{center}
\clearpage

\sect{8. Самостоятельная тренировка --- Путь А}
\textbf{10 заданий.} Выполняйте задания последовательно. К задачам на доказательство обязательно запишите обоснование; в задачах на графики укажите характерные точки и исключения.
\begin{enumerate}[label=\textbf{\arabic*.}]
\item Постройте график $y=-|x^2-4x+3|$. Укажите его нули и наибольшее значение.
\item В трапеции $ABCD$ основания $AD\parallel BC$. Биссектрисы углов при боковой стороне $AB$ пересекаются в точке $K$. Найдите $AB$, если $AK=9$, $BK=12$.
\item В трапеции $ABCD$ основания $AD\parallel BC$, $\angle DAB=68^\circ$. Биссектрисы углов $DAB$ и $ABC$ пересекаются в точке $K$. Найдите $\angle ABK$ и $\angle AKB$.
\item В тупоугольном треугольнике $ABC$ угол $A$ тупой. Высоты $BB_1$ и $CC_1$ проведены к прямым $AC$ и $AB$. Докажите, что $\triangle AB_1C_1\sim\triangle ABC$.
\item Решите неравенство $\displaystyle\frac{-8}{(x-4)^2-9}\ge0$.
\item Решите неравенство $\displaystyle\frac{x-3}{(x+1)(x-4)}>0$ методом интервалов.
\item В двух сосудах находятся $30$ кг и $20$ кг растворов кислоты. Если смешать их целиком, получится $40\%$-й раствор. Если смешать равные массы исходных растворов, получится $45\%$-й раствор. Найдите процент кислоты в первом сосуде.
\item Упростите выражение и постройте график
\[
h(x)=\frac{2|x|-3}{|x|-\frac23x^2}.
\]
Укажите область определения, координаты выколотых точек и вертикальную асимптоту.
\item Для функции $h(x)$ из задания 8 найдите все значения $k$, при которых прямая $y=kx$ не имеет с графиком общих точек.
\item Биссектрисы углов $A$ и $D$ трапеции $ABCD$ пересекаются в точке $M$, лежащей на стороне $BC$. Докажите, что расстояния от $M$ до прямых $AB$, $AD$ и $CD$ равны.
\end{enumerate}
\clearpage

\sect{9. Ответы и ориентиры для самопроверки}
\begin{enumerate}[label=\textbf{\arabic*.},itemsep=6pt]
\item Нули $x=1$, $x=3$; наибольшее значение $0$. На $[1;3]$ график задаётся $y=x^2-4x+3$, вне этого отрезка --- $y=-x^2+4x-3$.
\item $AB=15$: из $\angle AKB=90^\circ$ и теоремы Пифагора.
\item $\angle ABC=112^\circ$, поэтому $\angle ABK=56^\circ$ и $\angle AKB=90^\circ$.
\item Четыре точки $B,B_1,C,C_1$ лежат на окружности с диаметром $BC$. Отсюда $\angle AB_1C_1=\angle ABC$. Также $\angle B_1AC_1=\angle BAC$. Подобие доказано по двум углам.
\item $x\in(1;7)$; границы исключены, поскольку знаменатель обращается в нуль.
\item $x\in(-1;3)\cup(4;+\infty)$.
\item $20\%$ в первом сосуде, $70\%$ во втором. Проверка: $30\cdot0{,}2+20\cdot0{,}7=20$ кг кислоты из $50$ кг раствора.
\item $h(x)=-3/|x|$ при $x\ne0,\ \pm\frac32$. Выколоты точки $(-\frac32;-2)$ и $(\frac32;-2)$; вертикальная асимптота $x=0$.
\item $k\in\left\{-\frac43,\ 0,\ \frac43\right\}$. Ненулевые значения соответствуют прямым через удалённые точки графика.
\item Так как $M$ лежит на биссектрисе угла $A$, то $d(M,AB)=d(M,AD)$. По биссектрисе угла $D$: $d(M,AD)=d(M,CD)$. Следовательно, все три расстояния равны.
\end{enumerate}
\vfill
{\small\color{softgray}Самопроверка: при решении дробей сначала отмечайте ОДЗ; при доказательствах объясняйте равенство углов; в задачах на смеси проверяйте массовый баланс.}
\end{document}
